\section{C 异常处理：\texttt{isr\_handler\_c}}

\codefile{src/kernel/arch/idt.c}
\begin{lstlisting}[style=cstyle]
void isr_handler_c(const regs_t* r) {
    // Page Fault 特殊路径
    if (r->int_no == 14) {
        page_fault_handler(r);
        printk("Kernel panic: #PF Page Fault\n");
        for (;;) __asm__ volatile("hlt");
    }

    // 打印异常信息
    printk("\n[isr] %s vector=%u err=0x%08x eip=0x%08x\n",
           exception_name(r->int_no),
           r->int_no, r->err_code, r->eip);

    // Breakpoint (int3) 允许返回 → 用于自检
    if (r->int_no == 3) return;

    // 其他异常不可恢复
    printk("Kernel panic: %s\n", exception_name(r->int_no));
    for (;;) __asm__ volatile("hlt");
}
\end{lstlisting}

\subsection{Page Fault Handler：\reg{CR2} 与错误码}

\codefile{src/include/arch/idt.h}
\begin{lstlisting}[style=cstyle]
static inline uint32_t read_cr2(void) {
    uint32_t val;
    __asm__ volatile("mov %%cr2, %0" : "=r"(val));
    return val;
}
\end{lstlisting}

\reg{CR2} 在 \#PF 发生时由 CPU 自动写入故障虚拟地址。

错误码的逐 bit 含义：

\begin{figure}[H]
  \centering
  % figures/ch04/fig06-pf-5.tex — auto-extracted from chapter source
\begin{tikzpicture}
  \foreach \i/\label in {0/P, 1/W/R, 2/U/S, 3/RSVD, 4/I/D} {
    \node[bitcell, fill=white, minimum width=1.2cm] (pf\i) at (\i*1.3, 0) {\label};
    \node[font=\tiny\ttfamily, above=0.1cm of pf\i] {bit \i};
  }
\end{tikzpicture}

  \caption{\#PF 错误码的低 5 位}
\end{figure}

\begin{bitfield}
  \bit{0}{P (Present)：0=页不存在（not-present），1=权限违规（页存在但不允许访问）}
  \bit{1}{W/R：0=读操作引起，1=写操作引起}
  \bit{2}{U/S：0=在 Ring 0（内核态）发生，1=在 Ring 3（用户态）发生}
  \bit{3}{RSVD：PTE 中的保留位被错误置位（硬件级错误，通常是写 PTE 的代码有 bug）}
  \bit{4}{I/D：1=取指令（instruction fetch）引起，0=数据访问引起}
\end{bitfield}

常见错误码示例：

\begin{longtable}{rl}
  \toprule
  \textbf{err\_code} & \textbf{含义} \\
  \midrule
  \hex{00} & 内核态读一个不存在的页 \\
  \hex{02} & 内核态写一个不存在的页 \\
  \hex{05} & 用户态读一个存在但权限不足的页 \\
  \hex{07} & 用户态写一个存在但权限不足的页 \\
  \bottomrule
\end{longtable}

% ============================================================================

